\editionPage[\pdfbookmark[0]{Automatismes : questions flash}{bm-automatismes-1-parent}]{automatismes-1}{Automatismes · Questions flash}{Séries A et B · Nombres, calculs et proportions}{Sans calculatrice · \hyperref[automatismes-cor-1]{Corrigé p.\,\pageref*{automatismes-cor-1}}}
\begin{blocEdition}{automatismes-1}
Ces questions flash entretiennent les automatismes de la partie transversale du programme : calcul, proportions, évolutions, fonctions, géométrie, statistiques et probabilités. Une série se traite sans calculatrice, à l’oral ou à l’écrit, en 5 à 10 minutes ; la reprendre plusieurs fois dans l’année, en commençant dès les premiers chapitres. Les corrigés indiquent le bloc du programme de chaque question.
\editionEtape{Série A · Nombres, calculs et grandeurs}
\begin{deuxcolonnes}\begin{questions}
\item \begin{minipage}[t]{\linewidth}\raggedright\leavevmode Calculer $\dfrac23+\dfrac34$, puis comparer $\dfrac58$ et $\dfrac35$.\end{minipage}
\item \begin{minipage}[t]{\linewidth}\raggedright\leavevmode Écrire $\dfrac{2^5\times2^{-2}}{2^4}$ sous la forme $2^n$.\end{minipage}
\item \begin{minipage}[t]{\linewidth}\raggedright\leavevmode Écrire $0{,}35$ sous forme de fraction irréductible, puis de pourcentage.\end{minipage}
\item \begin{minipage}[t]{\linewidth}\raggedright\leavevmode Donner un ordre de grandeur de $49{,}7\times0{,}21$.\end{minipage}
\item \begin{minipage}[t]{\linewidth}\raggedright\leavevmode Comparer $a=\dfrac{1001}{1000}$ et $b=\dfrac{1000}{999}$ par leur quotient.\end{minipage}
\item \begin{minipage}[t]{\linewidth}\raggedright\leavevmode Convertir 90~km/h en m/s, puis $2{,}5$~m$^2$ en cm$^2$.\end{minipage}
\item \begin{minipage}[t]{\linewidth}\raggedright\leavevmode Simplifier $-(3-x)$ et $\dfrac2{3/4}$.\end{minipage}
\item \begin{minipage}[t]{\linewidth}\raggedright\leavevmode « Diminuer de 8\,\% » revient à multiplier par quel nombre ? Que signifie « multiplier par 1,15 » ?\end{minipage}
\item \begin{minipage}[t]{\linewidth}\raggedright\leavevmode Un calcul de probabilité donne $1{,}2$. Que conclure ?\end{minipage}
\item \begin{minipage}[t]{\linewidth}\raggedright\leavevmode $ABC$ est rectangle en $A$, avec $AB=4$ et $BC=8$. Calculer $\cos\widehat{ABC}$ et en déduire l’angle.\end{minipage}
\end{questions}\end{deuxcolonnes}
\editionEtape{Série B · Calcul littéral, équations et proportions}
\begin{deuxcolonnes}\begin{questions}
\item \begin{minipage}[t]{\linewidth}\raggedright\leavevmode Développer $(2x+3)^2$.\end{minipage}
\item \begin{minipage}[t]{\linewidth}\raggedright\leavevmode Factoriser $5x^2-15x$, puis $x^2-49$.\end{minipage}
\item \begin{minipage}[t]{\linewidth}\raggedright\leavevmode Résoudre $x^2=36$.\end{minipage}
\item \begin{minipage}[t]{\linewidth}\raggedright\leavevmode Résoudre $3x-4=x+6$.\end{minipage}
\item \begin{minipage}[t]{\linewidth}\raggedright\leavevmode Résoudre $\dfrac6x=2$ pour $x\ne0$.\end{minipage}
\item \begin{minipage}[t]{\linewidth}\raggedright\leavevmode Résoudre $-2x+5>1$.\end{minipage}
\item \begin{minipage}[t]{\linewidth}\raggedright\leavevmode Isoler $R$ dans $U=RI$ ($I\ne0$), puis $h$ dans $V=\dfrac13Bh$ ($B>0$).\end{minipage}
\item \begin{minipage}[t]{\linewidth}\raggedright\leavevmode Calculer $E=\dfrac12mv^2$ pour $m=2$~kg et $v=3$~m/s.\end{minipage}
\item \begin{minipage}[t]{\linewidth}\raggedright\leavevmode 9 élèves représentent 30\,\% d’une classe. Combien d’élèves compte la classe ?\end{minipage}
\item \begin{minipage}[t]{\linewidth}\raggedright\leavevmode Calculer 15\,\% de 240, puis écrire 15\,\% sous forme décimale et fractionnaire.\end{minipage}
\end{questions}\end{deuxcolonnes}

\end{blocEdition}
\editionPage{automatismes-2}{Automatismes · Questions flash}{Séries C et D · Géométrie, fonctions, données et hasard}{Sans calculatrice · \hyperref[automatismes-cor-2]{Corrigé p.\,\pageref*{automatismes-cor-2}}}
\begin{blocEdition}{automatismes-2}
\editionEtape{Série C · Géométrie et mesures}
\begin{deuxcolonnes}\begin{questions}
\item \begin{minipage}[t]{\linewidth}\raggedright\leavevmode Lire les abscisses des points $A$ et $B$.\par\smallskip\begin{tikzpicture}[x=1.75cm,font=\small]
\draw[otMuted,->](-2.3,0)--(1.35,0);
\foreach \k in {-8,...,4}{\draw[otMuted](\k/4,-.06)--(\k/4,.06);}
\foreach \x in {-2,-1,0,1}{\draw[otMuted](\x,-.11)--(\x,.11);\node[below=1.5mm,inner sep=1pt]at(\x,0){$\x$};}
\fill[otNavy](-1.25,0)circle(1.8pt)node[above=1mm]{$A$};\fill[otNavy](.75,0)circle(1.8pt)node[above=1mm]{$B$};
\end{tikzpicture}\end{minipage}
\item \begin{minipage}[t]{\linewidth}\raggedright\leavevmode Périmètre et aire d’un disque de rayon 3~cm.\end{minipage}
\item \begin{minipage}[t]{\linewidth}\raggedright\leavevmode Aire d’un triangle de base 7~cm et de hauteur 4~cm.\end{minipage}
\item \begin{minipage}[t]{\linewidth}\raggedright\leavevmode Volume d’un pavé de 2~cm sur 3~cm sur 5~cm ; d’une pyramide de même base et même hauteur.\end{minipage}
\item \begin{minipage}[t]{\linewidth}\raggedright\leavevmode Volume d’un cylindre, puis d’un cône, de rayon 2~cm et de hauteur 5~cm.\end{minipage}
\item \begin{minipage}[t]{\linewidth}\raggedright\leavevmode Volume d’un prisme droit de base 6~cm$^2$ et de hauteur 10~cm ; d’une boule de rayon 3~cm.\end{minipage}
\item \begin{minipage}[t]{\linewidth}\raggedright\leavevmode Les côtés de l’angle droit d’un triangle rectangle mesurent 6~cm et 8~cm. Hypoténuse ?\end{minipage}
\item \begin{minipage}[t]{\linewidth}\raggedright\leavevmode $M\in[AB]$, $N\in[AC]$, $(MN)\parallel(BC)$, $AM=2$, $AB=6$, $BC=9$. Calculer $MN$.\end{minipage}
\item \begin{minipage}[t]{\linewidth}\raggedright\leavevmode Triangle rectangle d’hypoténuse 10 et d’angle aigu $30^\circ$ : longueur du côté opposé ?\end{minipage}
\item \begin{minipage}[t]{\linewidth}\raggedright\leavevmode Convertir $1{,}5$~L en cm$^3$, puis 2~h~15~min en heures.\end{minipage}
\end{questions}\end{deuxcolonnes}
\editionEtape{Série D · Fonctions, statistiques et probabilités}
\begin{deuxcolonnes}\begin{center}\begin{tikzpicture}[x=.62cm,y=.42cm,font=\small]
\draw[otRule!70,step=1](-3,-3)grid(5,4);\draw[otMuted,->](-3,0)--(5.3,0)node[right]{$x$};\draw[otMuted,->](0,-3)--(0,4.3)node[above]{$y$};
\foreach \x in {-2,-1,0,1,2,3,4}{\node[below,inner sep=1.5pt]at(\x,-3){$\x$};}\foreach \y in {-2,-1,0,1,2,3}{\node[left,inner sep=1.5pt]at(-3,\y){$\y$};}
\draw[geoTrait](-2,3)--(0,-1)--(4,3);\node[otNavy,above right]at(4,3){$\mathcal C_f$};
\fill[otGreenDark](3,-2)circle(2pt)node[right]{$M$};
\end{tikzpicture}\par\legende{Courbe de $f$ formée de deux segments, et point $M$. Repère orthogonal, graduations lues sur les bords.}\end{center}
\begin{questions}
\item \begin{minipage}[t]{\linewidth}\raggedright\leavevmode Lire sur la courbe l’image de 0 par $f$.\end{minipage}
\item \begin{minipage}[t]{\linewidth}\raggedright\leavevmode Lire les antécédents de 1 par $f$.\end{minipage}
\item \begin{minipage}[t]{\linewidth}\raggedright\leavevmode Lire les coordonnées du point $M$.\end{minipage}
\item \begin{minipage}[t]{\linewidth}\raggedright\leavevmode $A(2;7)$ est-il sur la courbe d’équation $y=x^2+3$ ? Ordonnée du point d’abscisse $-1$ ?\end{minipage}
\item \begin{minipage}[t]{\linewidth}\raggedright\leavevmode Parmi $f(x)=3x$, $g(x)=2x-1$ et $h(x)=x^2$, lesquelles sont affines ? linéaires ?\end{minipage}
\item \begin{minipage}[t]{\linewidth}\raggedright\leavevmode Dans un diagramme semi-circulaire, quelle proportion représente un secteur de $45^\circ$ ?\end{minipage}
\item \begin{minipage}[t]{\linewidth}\raggedright\leavevmode Moyenne et médiane de la série $3$ ; $5$ ; $6$ ; $10$.\end{minipage}
\item \begin{minipage}[t]{\linewidth}\raggedright\leavevmode Classes $[0;10[$ d’effectif 3 et $[10;20[$ d’effectif 7 : estimer la moyenne avec les centres.\end{minipage}
\item \begin{minipage}[t]{\linewidth}\raggedright\leavevmode Boîtes à moustaches (min ; $Q_1$ ; médiane ; $Q_3$ ; max) : $A=(2;6;9;12;15)$ et $B=(4;7;9;10;14)$. Comparer.\end{minipage}
\item \begin{minipage}[t]{\linewidth}\raggedright\leavevmode On tire au hasard une carte parmi dix cartes numérotées de 1 à 10. Probabilité d’un multiple de 3 ? de l’événement contraire ?\end{minipage}
\end{questions}\end{deuxcolonnes}

\end{blocEdition}
